\documentclass[11pt]{article}
\usepackage[T1]{fontenc}
\usepackage{amsmath,amssymb,amsthm}
\usepackage[hidelinks]{hyperref}
\newcommand{\NN}{\mathbb{N}}
\newtheorem{theorem}{Theorem}[section]
\theoremstyle{definition}
\newtheorem{definition}{Definition}[section]
\title{A Synthetic Example: Sums of Consecutive Odd Integers}
\author{Example Author}
\date{}
\hypersetup{pdfauthor={Example Author},pdftitle={Synthetic MPS Demo},pdfsubject={Fictional elementary writing example}}
\begin{document}
\maketitle
\begin{abstract}
This synthetic writing example proves the classical identity expressing a square
as a sum of consecutive odd integers. It demonstrates a small LaTeX project and
does not claim a new mathematical result.
\end{abstract}
\tableofcontents
\section{Introduction}
Our example is the identity in Theorem~\ref{thm:odd}. All names and project
details are placeholders. No external research report is used.
\section{Preliminaries}
\begin{definition}
For $n\in\NN=\{1,2,3,\ldots\}$, set $S_n=\sum_{k=1}^{n}(2k-1)$.
\end{definition}
\section{A Telescoping Calculation}
The identity $k^2-(k-1)^2=2k-1$ expresses each odd integer as the difference
between consecutive squares. Summing these differences cancels all interior
terms.
\section{Main Result}
\begin{theorem}\label{thm:odd}
For every $n\in\NN$, one has $S_n=n^2$.
\end{theorem}
\begin{proof}
Using the preceding calculation gives
\[
 S_n=\sum_{k=1}^{n}\bigl(k^2-(k-1)^2\bigr)=n^2-0^2=n^2.
\]
This holds for every positive integer $n$.
\end{proof}
\section{An Application}
For $n=4$, the formula gives $1+3+5+7=16$. More generally,
$S_{n+1}-S_n=2n+1$ for $n\in\NN$. This application illustrates the statement;
it is not an independent computational verification service.
\end{document}
